\section{Introduction}
\label{sec:intro}

If you export every song you have ever streamed and paste the log into a frontier LLM, will the resulting recommendations be better than the ones the platform serves you? This is the consumer-facing form of an open empirical question in the recommender systems community: how much of the engineering investment in collaborative and content-based filtering is replaceable by prompting a single foundation model? The answer matters in two directions. If LLMs are competitive, users gain a path to platform-independent recommendation and researchers must rethink which classical techniques retain value. If they are not, the result tempers a wave of LLM enthusiasm with a clean, controlled comparison.

\para{Why music, and why now.} Music recommendation is the right testbed: histories are dense and time-stamped, the catalog is finite, and standard offline metrics are stable. Frontier LLMs are now cheap and fast enough that a per-user inference call costs roughly one cent. Yet the literature lacks a head-to-head comparison that simultaneously measures rank quality, popularity bias, hallucination, and personalization gain on real listening histories \citep{epure2025music}. The closest precedent, \citet{boadana2025llm}, compares a multi-agent LLM only against a TF-IDF content-based filter on $19$ Spotify-export users and reports subjective like-rate, not held-out accuracy. \citet{hou2024llm} establish that LLMs can rank but evaluate on movies and books, not music. \citet{sguerra2025biases} show that user-perceived profile quality and downstream recall are weakly correlated, so subjective gains do not transfer.

\para{What we do.} We restate the user's question as four falsifiable hypotheses (relevance, no-popularity-shortcut, low-hallucination, real-personalization) and test them with a single experimental protocol: every method ranks the \emph{same} $20$-track candidate pool, so differences are attributable to the ranking decision rather than candidate generation. We use \claude with a Hou-et-al.~zero-shot ranking prompt, fed the user's last $50$ tracks from a Last.fm 1K~\citep{schedl2017lfm1b} export. We compare against four classical baselines that operationally proxy what a Spotify-style stack does: popularity, item-kNN collaborative filtering, TF-IDF content-based filtering, and random. We add \gpt as a cross-model robustness check and a no-history ablation to isolate personalization from priors.

\para{What we find.} On $N{=}24$ users with real $30$-day held-out windows, \claude with history achieves the highest \ndcg ($0.375$), tying \gpt ($0.376$) and beating popularity by $+0.207$ \ndcg ($p{=}0.028$, Cohen's $d{=}0.53$). Against the strongest classical baseline (\cbf, \ndcg $0.333$) the LLM is nominally better ($+0.042$) but not significantly so at this sample size, while it does not lose on any metric. \claude does \emph{not} win by exploiting popularity: its \arp ($58.6$) is lower than \cbf's ($59.7$). Removing history collapses \claude to \ndcg $0.082$ -- below random -- and the personalization gain is large ($+0.29$ \ndcg, $p{=}0.002$, $d{=}0.86$), confirming that the LLM is using the listening log rather than relying on priors. Free-generation hallucination is high in absolute terms ($50\%$ of recommendations resolve to our $77$K-track catalog), but spot-checking shows most unresolved recommendations are real songs absent from this catalog snapshot.

\para{Contributions.} Our contributions are:
\begin{itemize}[leftmargin=*,itemsep=0pt,topsep=0pt]
    \item We propose a clean head-to-head protocol that ranks an identical candidate pool with $7$ methods, isolating the ranking decision from candidate generation, and that scores all methods on a four-axis evaluation (relevance, popularity bias, hallucination, personalization gain) following \citet{epure2025music}.
    \item We conduct paired-Wilcoxon comparisons across two frontier LLMs (\claude and \gpt) and four classical baselines on $N{=}24$ Last.fm users, finding that LLM ranking matches or exceeds every baseline on rank-position-sensitive metrics while spending only $\$0.45$ in API calls.
    \item We complement the ranking experiment with a no-history ablation showing a large personalization gain ($d{=}0.86$, $p{=}0.002$), establishing that the LLM is conditioning on the history rather than reproducing popular-track priors.
    \item We release the end-to-end pipeline, including pool construction, prompts, raw LLM outputs, and per-user statistics, to make Last.fm-as-Spotify-export a reproducible benchmark for future LLM-recommender studies.
\end{itemize}

\para{Paper organization.} \secref{sec:related} positions our work against the LLM-recommender literature. \secref{sec:method} describes the data, candidate pool, methods, and metrics. \secref{sec:results} reports the headline ranking comparison, popularity-bias diagnostic, hallucination measurement, personalization-gain ablation, and cost/latency accounting. \secref{sec:discussion} discusses limitations and broader implications, and \secref{sec:conclusion} concludes.
